% ===================================================================
% chapters/02a_literature_target.tex -- writer L1
% Literature basis of the TCAD model: opening, source map, citation
% conventions, and the page-exact treatment of the target paper
% (Januar et al. 2026, Small Structures).
% ASCII only. Page numbers follow CONVENTIONS section 5.
% ===================================================================
\chapter{Literature basis of the TCAD model}\label{ch:literature}

The credibility of a drift--diffusion model of an ultrathin oxide-semiconductor transistor depends
on the numbers that enter it. The thickness-dependent indium--tungsten oxide (\IWO) model of this
report contains about forty inputs: a gate stack, a band structure that changes with film thickness,
an electron effective mass, a trap density of states, fixed and interface charges, a band mobility,
contact assumptions and numerical settings. Some of these inputs were measured on the target device,
some were taken from measurements on a related material (pure \InO{}) and some from
density-functional (DFT) slab calculations on pure \InO{}; others were assumed, and the remainder
were fitted to the measured transfer curves. For every source that supplied a number or a modelling
choice, this chapter addresses three questions: (i) what exactly was taken, together with the printed
page (and figure or equation) at which it can be checked; (ii) how the item enters the ATLAS model,
with a reference to the parameter section in which it is derived; and (iii) what was considered but
deliberately not used, and where the sources disagree. The chapter closes with a master provenance
table (\cref{sec:lit-provenance}), which lists every model input with its status label, its source
page or the run in which it was fitted, and the runs that test its sensitivity.

\Cref{tab:lit-sources} gives an overview of all sources. Each bundled paper (text available locally
with page markers) is treated in a section of its own, in
\cref{sec:lit-januar,sec:lit-si2021,sec:lit-lin2022,sec:lit-si2022,sec:lit-wang2022,sec:lit-kim2024,sec:lit-stokey2021,sec:lit-pandey2025,sec:lit-wangx2025,sec:lit-anusha2024,sec:lit-giannozzi2009}.
The sources that are not bundled appear only through the bundled record that cites them
(\cref{sec:lit-secondary}), and the Silvaco manual pages behind every deck statement are listed in
\cref{sec:lit-atlas-manual}.

\begin{table}[htbp]
\centering
\footnotesize
\caption{Map of all literature and reference sources used by the model. ``Bundled'' means that the
text is available locally with page markers and every page cited in this chapter was located in it.
Status labels follow \cref{sec:lit-convention}.}
\label{tab:lit-sources}
\begin{tabularx}{\linewidth}{@{}>{\raggedright\arraybackslash}p{2.0cm}>{\raggedright\arraybackslash}p{3.3cm}>{\raggedright\arraybackslash}p{1.0cm}>{\raggedright\arraybackslash}X>{\raggedright\arraybackslash}p{1.5cm}@{}}
\toprule
Key & Work & Bundled & Role in the model & Section \\
\midrule
Januar2026 & Target paper: \IWO{} and \InO{} TFTs vs thickness, Small Struct.\ & yes & gate stack and process; roughness $\Dsr$; trap trend $\Nt(t)$; mobility extraction convention; analytic framework & \cref{sec:lit-januar} \\
JanuarSI2026 & Supporting information of the target paper & no & permittivities of \IWO, \InO, \HfO{} (Fig.~S1b); grains; temperature data (Fig.~S12) & \cref{sec:lit-secondary} \\
Si2021 & \InO{} transistors down to \SI{0.7}{nm}, Nano Lett.\ & yes & confinement point $\dEc(\SI{1.5}{nm})$; $\Ec$-only partition; CNL above $\Ec$ & \cref{sec:lit-si2021} \\
Lin2022 & nm-thick \InO{} with high drain current, ACS Nano & yes & PBE slab gaps and masses; $\Dit$ of \HfO/\InO & \cref{sec:lit-lin2022} \\
Si2022 & Scaled ALD \InO{} transistors, Nat.\ Electron.\ & yes & CNL and low contact resistance (Ohmic S/D, qualitative) & \cref{sec:lit-si2022} \\
Wang2022 & Traps in \HfO/\InO{} MOS capacitors, Front.\ Mater.\ & yes & $\Dit$ bracket; subgap DOS and TCAD trap conventions & \cref{sec:lit-wang2022} \\
Kim2024 & Thickness-dependent stability of \IWO{} TFTs, APL & yes & thickness trend of bulk/border traps (support for $\Ndeff(t)$ rise) & \cref{sec:lit-kim2024} \\
Stokey2021 & Dielectric constants and electron mass of \InO, JAP & yes & bulk $\mstar = 0.208\,m_0$; $\varepsilon_{\mathrm{DC}} = 10.55$ (sensitivity bound) & \cref{sec:lit-stokey2021} \\
Pandey2025 & Analytical short-channel model of oxide TFTs, TED & yes & method reference only & \cref{sec:lit-pandey2025} \\
WangX2025 & Physics-based compact model of IGZO FETs, TED & yes & method reference (Gaussian shallow donors) & \cref{sec:lit-wangx2025} \\
Anusha2024 & Review of IGZO TCAD/compact modelling, Meas.\ Sens.\ & yes & SS--trap relation (theory) & \cref{sec:lit-anusha2024} \\
Giannozzi2009 & Quantum ESPRESSO, JPCM & yes & software reference of the unexecuted DFT workflow & \cref{sec:lit-giannozzi2009} \\
Fan2021, Hamberg1986, Yoo2024, Lee2018, King2008, Robertson2011, Hinuma2019, Ando1982 & secondary sources & no & only through a citing bundled record or the Astra audit & \cref{sec:lit-secondary} \\
AtlasManual, DeckBuildManual & Silvaco ATLAS 5.28.1.R / DeckBuild manuals & local PDF & syntax and physics of every deck statement & \cref{sec:lit-atlas-manual} \\
ExperimentalWorkbook, DeviceSchematic & measured ID--VG workbook; device schematic & local & ground truth; geometry & \cref{ch:device} \\
\bottomrule
\end{tabularx}
\end{table}

\section{Citation and page conventions}\label{sec:lit-convention}

\paragraph{Page-exact citations.}
Every statement taken from a paper carries the printed page of the article, written as
\citep[p.~21539]{Lin2022}; figures and equations are added where useful, as in
\citep[p.~8, Fig.~5f]{Januar2026}. The local text files contain one block per PDF page. Because the
mapping from the PDF page to the printed article page differs between journals, it is fixed in
\cref{tab:lit-pages}. Where a journal uses article-number pagination (APL, JAP), the page is
written as \emph{article number--page}, as in \citep[p.~173507-6]{Kim2024}. The publisher cover
pages of Kim2024 and Stokey2021 carry no article page and are never cited.

\begin{table}[htbp]
\centering
\footnotesize
\caption{Page convention: printed page to be cited for the $N$-th PDF page of each bundled source.}
\label{tab:lit-pages}
\begin{tabularx}{\linewidth}{@{}>{\raggedright\arraybackslash}p{2.1cm}>{\raggedright\arraybackslash}X>{\raggedright\arraybackslash}p{4.2cm}>{\raggedright\arraybackslash}p{2.6cm}@{}}
\toprule
Key & Local text file & Printed page for PDF page $N$ & Printed range \\
\midrule
Januar2026 & \file{Januar2026_SmallStruct_IWO.txt} & $N$ (``$N$ of 12'') & 1--12 \\
Si2021 & \file{Si2021_NanoLett_In2O3_0p7nm.txt} & $499 + N$ & 500--506 \\
Lin2022 & \file{Lin2022_ACSNano_thick_In2O3.txt} & $21535 + N$ & 21536--21545 \\
Si2022 & \file{NatElec2022_s41928-022-00718-w.txt} & $163 + N$ & 164--170 \\
Wang2022 & \file{Wang2022_FrontMater_traps_HfO2In2O3.txt} & $N$ & 1--7 \\
Kim2024 & \file{Kim2024_APL_IWO_thickness_stability.txt} & 173507-$(N-1)$; PDF p.~1 is a cover & 173507-1 to -7 \\
Stokey2021 & \file{Stokey2021_JAP_In2O3_dielectric_mass.txt} & 225102-$(N-1)$; PDF p.~1 is a cover & 225102-1 to -15 \\
Pandey2025 & \file{Pandey2025_TED_SCE_TFT.txt} & $N$ (early-access pagination) & 1--9 \\
WangX2025 & \file{Wang2025_TED_IGZO_compact.txt} & $2389 + N$ & 2390--2398 \\
Anusha2024 & \file{Anusha2024_MeasSens_IGZO_review.txt} & $N$ & 1--16 \\
Giannozzi2009 & \file{Giannozzi2009_QE.txt} & $N$ & 1--19 \\
AtlasManual & \file{atlas_users1.pdf} (ATLAS 5.28.1.R) & printed footer page; PDF page stated where they differ & see \cref{tab:lit-atlas} \\
\bottomrule
\end{tabularx}
\end{table}

\paragraph{Sources that are not bundled.}
The supporting information of the target paper and eight secondary works (\cref{sec:lit-secondary})
are not available locally, and they are never cited as if they had been read. Each of them appears
either together with the bundled record that cites it and that record's page (``as cited in
\citet[p.~501]{Si2021}''), or with an explicit note that the value was taken over from the audit of
the earlier Astra package
(\file{inputs/astra_IWO_PRIORITY_THREE_20260912/.../audit/docs/paper_revision_20260911/}).

\paragraph{Status labels.}
Every model input carries one label (typeset as small capitals, e.g.\ \prov{FITTED}):
\prov{MEASURED\_TARGET\_DEVICE} (measured on the target process or device);
\prov{LITERATURE\_IWO} (reported for \IWO, possibly a different process or device);
\prov{MEASURED\_RELATED\_MATERIAL} (measured on pure \InO{});
\prov{LITERATURE\_IN2O3\_PROXY} (reported for \InO{} and transferred as a proxy);
\prov{DFT\_PURE\_IN2O3\_PROXY} (PBE slab-minus-bulk differences for pure \InO{});
\prov{CALCULATED} (computed from other inputs by a stated formula);
\prov{ASSUMED} (chosen without a direct measurement; a way to measure it is given in the parameter section);
\prov{FITTED} (adjusted to the measured transfer curves, with the run that fixed it).
The evidence classes of \cref{ch:validation} (\evDATA, \evNUM, \evCAL, \evOOS, \evPOST, \evPRED,
\evSURR, \evTHEORY, \evLIT, \evCORR) are used when a claim, rather than an input, is qualified. A
literature value used as a model input belongs to the class \evLIT; a fitted value belongs to the
class \evCAL{} and never counts as validation.

\paragraph{Structure of each paper section.}
Each paper section contains (a) a short summary; (b) a box that lists every item used, with the value,
the printed page, the way in which it enters the model, the parameter section of
\cref{ch:derivations-es} or \cref{ch:derivations-tr} in which it is derived, and its caveat; (c) the
items that were considered but not used; and (d) the contradictions with other sources or within the
paper.

\section{Januar et al.\ 2026: the target device and its analytic framework}\label{sec:lit-januar}

\paragraph{Summary.}
\citet{Januar2026} fabricate bottom-gate thin-film transistors with sputtered \IWO{} (``$\sim$2\% W'')
and \InO{} channels on a TiN/\HfO/\AlO{} gate stack with Pd source/drain contacts
\citep[pp.~2--3]{Januar2026}, and study how the channel thickness controls mobility, threshold,
subthreshold swing and bias stability. They report peak field-effect mobilities of about 5.1 to
$27.4\cmVs$ for the \IWO{} thickness series \citep[p.~3]{Januar2026}. The decline for thinner films
is interpreted within a trap-limited transport framework in which the tail-trap density $\Nt$ grows
roughly fourfold from 13.2 to \SI{2}{nm} \citep[p.~7]{Januar2026}, and the gate-dependent mobility
is described by a surface-roughness (Ando-type) term with an effective roughness
$\Dsr \approx \SI{2.87}{\angstrom}$ for \IWO{} \citep[pp.~7--8, Fig.~5f]{Januar2026}. PBE-GGA calculations
are used for qualitative trends in the band structure and at the \IWO/\HfO{} interface
\citep[p.~3]{Januar2026}, and positive-bias-stress (PBS) devices are analysed with the same compact
model \citep[p.~9]{Januar2026}. The measured transfer curves of the workbook used in this report
belong to the thickness-series devices of the paper; this identification was established in the S5
report, as described below.

\paragraph{Device and process facts used.}
The gate is a \SI{50}{nm} TiN electrode deposited by physical vapour deposition, followed by the
\HfO/\AlO{} gate dielectric and the \IWO{} or \InO{} channel \citep[p.~2]{Januar2026}. The channel
lacks long-range crystallinity and consists of nanoscale grains, which are smaller for \IWO{} than for
\InO{} at comparable thickness \citep[p.~2]{Januar2026}. Pd source/drain electrodes of \SI{70}{nm}
complete the device, which is then annealed by rapid thermal annealing \citep[p.~3]{Januar2026}. The
paper states that the trap metrics improve after annealing at \SI{150}{\degreeCelsius}, and that
higher annealing temperatures can raise the mobility further but make the device harder to turn off,
through a crystallinity-driven increase of the carrier concentration \citep[p.~3]{Januar2026}. The
\HfO{} and \AlO{} thicknesses (15 and \SI{2}{nm}), $W/L = 290/20$\,$\mu$m and $\Vd = \SI{0.7}{V}$ are
taken from the device schematic \citep{DeviceSchematic} and are consistent with the stack description
of the paper (details in \cref{sec:inp-device,sec:inp-geometry}).

\begin{usedbox}
\textbf{Items used from \citet{Januar2026} (and its SI, not bundled).}\par\smallskip
\footnotesize
\begin{tabularx}{\linewidth}{@{}>{\raggedright\arraybackslash}p{1.9cm}>{\raggedright\arraybackslash}p{1.8cm}>{\raggedright\arraybackslash}p{1.5cm}>{\raggedright\arraybackslash}X>{\raggedright\arraybackslash}p{1.35cm}>{\raggedright\arraybackslash}X@{}}
\toprule
Quantity & Value & Page & How it enters the model & Section & Caveat \\
\midrule
Gate metal & TiN, \SI{50}{nm} & p.~2 & gate electrode region; work function assumed \SI{4.70}{eV} & \cref{sec:par-stack,sec:par-gatewf} & TiN work function not reported \\
Gate dielectric & \HfO/\AlO{} bilayer & p.~2 & two insulator regions; $\Cox = \SI{8.955e-7}{\farad\per\centi\metre\squared}$ & \cref{sec:par-stack} & thicknesses from the schematic \\
Channel & \IWO{} ($\sim$2\% W) & p.~2 & material label only & \cref{sec:par-nd} & composition basis not defined; never converted to a donor density \\
S/D metal & Pd, \SI{70}{nm} & p.~3 & ideal Ohmic contacts & \cref{sec:par-contacts} & no TLM data \\
$\epsIWO$ & 9.30 & SI Fig.~S1b (caption list p.~11) & permittivity of the channel & \cref{sec:par-eps-iwo} & SI not bundled; value recorded by the Astra audit \\
$\varepsilon_{\mathrm{HfO_2}}$ & 19.57 & SI Fig.~S1b (p.~11) & gate-stack capacitance & \cref{sec:par-stack} & same \\
$\Dsr$ & \SI{2.87}{\angstrom} & p.~8, Fig.~5f & roughness factor $(1-\Dsr/t)^2 = 0.734$, 0.911, 0.957, 0.982 & \cref{sec:par-rough} & fitted by the authors on $\musat$ values; unit ambiguity between prose (\AA) and figure axis (nm) \\
$\Nt(13.2)/\Nt(2)$ & $\approx 1/4$ & p.~7, Fig.~5a & trend constraint for $\Nt(t) \propto t^{-0.75}$ & \cref{sec:par-tail} & a compact-model output, not a measurement \\
$\Nt$ convention & $n_t \approx \theta_t \Nt \exp[(\EF-\Ec)/\kB T_t]$ & pp.~4, 9 & mapping to ATLAS $\NTA\WTA$ & \cref{sec:par-tail,sec:der-tlc} & absolute $T_t$ not given (only changes) \\
$\mu_{\max}$ & 5.1 / $27.4\cmVs$ (2 / \SI{13.2}{nm}) & pp.~3, 7 & defines the extraction convention; reproduced by the saturation formula (5.09 / 27.40, S5 report) & \cref{sec:der-mufe} & not a model input; linear $\muFE$ is 12.15 / 50.17 \\
Roll-off $\mu_0/(1+\theta_r V_{\mathrm{ov}})$ & form only & p.~7 & NOT inserted; approximation gap & \cref{sec:par-mu,sec:der-mu-linear} & ATLAS ignores MOBILITY updates between SOLVEs \\
Smoother \IWO/\HfO{} interface & qualitative & pp.~2, 8 & supports the low assumed $\Dit$ & \cref{sec:par-dit} & the target interface is \IWO/\AlO \\
Grain/crystallinity trend & qualitative & pp.~2--3 & hypothesis for the $\muband$ step 6.3 $\rightarrow$ \SI{13.2}{nm} & \cref{sec:par-mu,sec:mech-transport} & not quantified \\
\bottomrule
\end{tabularx}
\end{usedbox}

\paragraph{How the paper's analytic framework is used.}
The paper bases its interpretation on a trap-limited conduction model in which the free-carrier
fraction of the induced charge sets the apparent mobility \citep[p.~7, Eq.~5]{Januar2026}, and on a
tail-trap density written with a factor $\theta_t$ and a characteristic temperature $T_t$
\citep[p.~4, Eq.~3]{Januar2026}; the
same convention is used for the PBS extraction \citep[p.~9]{Januar2026}. These closed forms are not
inserted into the TCAD model. The model solves the Poisson and drift--diffusion equations with an
explicit acceptor-like exponential tail $g_{\mathrm{TA}}(E) = \NTA\exp[(E-\Ec)/\WTA]$, whose integral
$\Nt \approx \NTA\WTA$ plays the role of $\Nt$ in the paper. The partition between free and trapped
charge then follows self-consistently from the occupancy of the discretised trap levels
(\cref{sec:der-tlc}). Only the \emph{trend} reported in the paper ($\times 4$ from 13.2 to
\SI{2}{nm}) constrains the model; the absolute magnitudes differ by the convention factor $\theta_t$
and are therefore compared but not equated (\cref{sec:par-tail}).

\paragraph{Mobility definition (resolved discrepancy).}
The mobilities of the paper (5.1 and $27.4\cmVs$; \citealp[pp.~3, 7]{Januar2026}) are not the
linear-regime field-effect mobilities of the workbook curves (12.15, 10.95 and $50.17\cmVs$
for 2.0, 6.3 and \SI{13.2}{nm}, \file{docs/MATERIAL_PARAMETER_EXTRACTION.md} section 10; the
$49.2\cmVs$ extracted at \SI{31.8}{nm} is not a mobility, because $\gm$ has several maxima there).
Applying the saturation formula $\mu = (2L/W\Cox)(\mathrm{d}\sqrt{\Id}/\mathrm{d}\Vg)^2_{\max}$ to
the workbook curves at $\Vd = \SI{0.7}{V}$ gives 5.09 and $27.40\cmVs$ (S5 report,
\file{analysis_2026-09-25/S5_contacts_experimental_validation.md}). This result identifies the
workbook curves with the thickness series of the paper, and the $\mu_{\max}$ of the paper with a
peak of the saturation formula. All comparisons in this report use one extraction, applied in the
same way to measured and simulated curves (\cref{sec:der-mufe,sec:der-metrics}).

\paragraph{Considered but not used.}
The absolute values $\Nt(\SI{2}{nm}) = \SI{5.3e19}{cm^{-3}}$ and
$\Nt(\SI{10}{nm}) = \SI{3.39e18}{cm^{-3}}$ of the unstressed PBS devices \citep[p.~9]{Januar2026}
were not used, because they belong to a different device set and depend on a $\theta_t$ convention;
the model anchor is instead the real-ATLAS tail integral $\NTA\WTA = \SI{2.0e19}{cm^{-3}}$ at
\SI{2}{nm} (\cref{sec:par-tail}). The stress-induced changes of $\Nt$ and $T_t$ under PBS
\citep[p.~9]{Januar2026} were also not used, since the model describes unstressed DC transfer curves
only. The PBE band structures and interface structures \citep[pp.~3, 5]{Januar2026} were used only
qualitatively (W raises $\mstar$; the \IWO/\HfO{} interface is smoother). No number was taken from
them, because absolute PBE values are not transferable and the effect of W on $\mstar$ cannot be
determined numerically from the available data. The oxygen-vacancy donor excess of undoped \InO{}
($n > 10^{19}\cmcu$; \citealp[p.~2]{Januar2026}) serves as the motivation for W doping and is not a
model input; the fitted $\Ndeff$ of \IWO{} is 2--3 orders of magnitude lower (\cref{sec:par-nd}).
The \InO{} reference devices of the paper consist of a different material and were not modelled.
Finally, the elevated-temperature transfer curves of the \SI{2}{nm} device at 338 and \SI{358}{K}
(SI Fig.~S12, referred to on p.~9 and captioned on p.~11) were not used for calibration; they form
the natural first test of the registered temperature predictions \run{0044}--\run{0047}
(\cref{sec:pred-temperature}).

\paragraph{Contradictions and open points.}
\begin{enumerate}
\item \emph{$\Nt(t)$ exponent.} The thickness-series ratio ($\approx 4$ from 13.2 to \SI{2}{nm},
  p.~7) implies $\Nt \propto t^{-0.73}$, whereas the two PBS devices ($5.3\times10^{19}$ at
  \SI{2}{nm}, $3.39\times10^{18}\cmcu$ at \SI{10}{nm}, p.~9) imply $t^{-1.71}$ (S2 report). The two
  come from different device sets; the model follows the thickness series (exponent 0.75).
\item \emph{Mobility extraction.} The $\mu_{\max}$ of the paper and the linear $\muFE$ differ by a
  factor of about 2 (see above); the difference arises from the extraction formula, not from the data.
\item \emph{Roughness units.} $\Dsr$ is printed as \SI{2.87}{\angstrom} in the text (p.~8), whereas the
  fitted figure uses a nm axis (a unit ambiguity recorded by the Astra audit); the \AA{} value is used.
\item \emph{Composition.} The specification ``$\sim$2\% W'' (p.~2) does not state whether an atomic,
  cation or WO$_3$ fraction is meant; it is therefore treated as a label only.
\item \emph{Equation conventions.} The Astra audit records that some main-text equations
  differ in convention from the SI. Only the $\Nt$ trend and the $\theta_t$ definition are used here,
  so the discrepancy does not propagate into the model.
\item \emph{Interface.} The paper discusses the \IWO/\HfO{} interface (pp.~2, 5, 8), but the target
  stack has a \SI{2}{nm} \AlO{} interlayer between \HfO{} and \IWO; statements about the interface are
  therefore transferred only qualitatively.
\end{enumerate}
